\Needspace{5\baselineskip}\pdfbookmark[2]{Zweite Abteilung}{witwen-II-2}
\begin{center}\textit{Zweite Abteilung.}\end{center}
\noindent\hspace*{\parindent}
Da ich die bei der Aufstellung der Bilanz der Witwenkasse einzuschlagenden Wege bereits in meinem Votum vom 8. Januar dieses Jahres ausführlich beschrieben und bei der Durchführung der Arbeit keinen Anlass zu Änderungen gefunden habe, kann ich mich auf einen Verweis darauf beschränken. Ich will hier nur bemerken, dass das Wesen der Methode darin besteht, den gegenwärtigen Geldwert der Verpflichtungen der Kasse unter den drei Rubriken zu ermitteln
\begin{quote}
Verpflichtungen gegenüber den jetzigen 19 Witwen,\par
Verpflichtungen gegenüber den Witwen (und Waisen) der jetzigen 51 Teilnehmer,\par
Verpflichtungen gegenüber den Witwen und Waisen der künftig Beitretenden,
\end{quote}
\noindent wobei für die zweite und dritte Rubrik der gegenwärtige Geldwert der Gegenleistungen durch die Beiträge abzuziehen ist. Zunächst schicke ich einige allgemeine Erläuterungen voraus.

Als Zeitpunkt, auf den alle künftigen Leistungen durch Diskontierung bezogen werden, ist der 1. Oktober 1845 gewählt. Dabei wird vorausgesetzt, dass die Witwen und Teilnehmer ihre bis dahin fällig gewordenen Pensionen beziehungsweise Beiträge bereits empfangen beziehungsweise geleistet haben.
\par\SourceLeaf{156}
Für die Rechnung habe ich die höchst wertvollen Sterblichkeitstafeln verwendet, die \textsc{Brune} aus den bei der preußischen Witwenkasse an 31500 Ehepaaren gemachten Erfahrungen abgeleitet hat. Nur für das Sterben der Männer im höchsten Lebensalter sind diese Tafeln mangelhaft, da die Unterlagen der Witwenkasse dafür keine hinreichenden Daten enthielten. Daher habe ich es vorgezogen, das Sterben der Männer über 80 Jahren nach denselben Verhältnissen zu berechnen, die die Tafeln für das weibliche Geschlecht angeben. Um die Veränderungen zu berechnen, die der Wert der Witwenpensionen durch die Berücksichtigung minderjähriger Kinder erfährt, sofern solche vorhanden sind, habe ich für diese die Sterblichkeitstafeln von \textsc{Deparcieux} verwendet. Im Übrigen hat die Wahl der Tafeln sowohl für die Sterblichkeit der Kinder als auch für die der Männer im höchsten Lebensalter nur einen sehr geringen Einfluss auf die Rechenergebnisse.

Alle Ergebnisse hängen wesentlich vom zugrunde gelegten Zinsfuß ab. Um hier nichts zu wünschen übrig zu lassen, habe ich daher sämtliche Rechnungen sowohl mit einem Zinsfuß von 4 Prozent als auch mit einem von \(3\tfrac{1}{2}\) Prozent durchgeführt, obwohl sich die Arbeit dadurch genau verdoppelte. Die Geldbeträge sind stets in Gold zu verstehen. Die Rechnung wurde durchgehend bis auf Bruchteile des Talers genau geführt; diese Bruchteile sind jedoch in der vorliegenden Abschrift weggelassen. Daraus kann sich bei den Summen hier und da ein Unterschied von einer oder ein paar Einheiten ergeben. Diesen geringfügigen Umstand bemerke ich hier nur, damit jemand, der etwa die eine oder andere Summe nachrechnet, solche Unterschiede nicht für Zeichen einer ungenauen Rechnung hält. Sie sind vielmehr gerade umgekehrt eine Folge der größeren Genauigkeit, die bei der Rechnung eingehalten wurde.

\Needspace{6\baselineskip}\pdfbookmark[3]{I. Berechnung des Geldwerts der Witwenpensionen}{witwen-II-2-I}
\par\medskip\noindent I.\quad Berechnung des auf den 1. Oktober 1845 zurückgerechneten Geldwerts der Witwenpensionen.\par\smallskip
Die Gesamtzahl aller seit Errichtung der Witwenkasse bis jetzt eingetretenen Pensionierungen beträgt 68. Um meine Arbeit möglichst leicht überschaubar zu machen, habe ich sie fortlaufend nummeriert. Die Ansätze in der folgenden Tabelle geben die Summen an, die den einzelnen Witwen ausgezahlt werden müssten, wenn ihre Ansprüche abgefunden werden sollten, oder umgekehrt die Summen, die die Witwen einzahlen müssten, wenn sie eine solche Berechtigung, wie sie ihnen jetzt zusteht, erst kaufen wollten. Der Rechnung liegt der jetzige Pensionssatz von 250 Rthl. zugrunde, mit Ausnahme der Nr. 58, bei der er 200 Rthl. beträgt. Die Witwen 50. 53. 58. 64. 68 haben Kinder unter 20 Jahren, die bei der Rechnung genau berücksichtigt sind, ebenso wie der Umstand, dass die Pensionen immer noch für den vollen Monat ausgezahlt werden, in dem ein Abgang erfolgt.

\begin{center}
\footnotesize\setlength{\tabcolsep}{3pt}
\begin{tabular}{r|c|r|r||r|c|r|r}
\multicolumn{2}{c|}{}&\multicolumn{2}{c||}{\shortstack{Gegenwärtiger Wert\\der Pension.}}&\multicolumn{2}{c|}{}&\multicolumn{2}{c}{\shortstack{Gegenwärtiger Wert\\der Pensionen.}}\\\cline{3-4}\cline{7-8}
Nr.& &\(3\tfrac{1}{2}\) Prozent&4 Prozent&Nr.& &\(3\tfrac{1}{2}\) Prozent&4 Prozent\\\hline
29&L.&957 Thl.&941 Thl.&58&G.&3172 Thl.&3000 Thl.\\
31&S.&1756 ,, &1708 ,, &59&W.&2850 ,, &2732 ,, \\
41&W.&1122 ,, &1101 ,, &60&S.&2606 ,, &2507 ,, \\
44&F.&2543 ,, &2448 ,, &62&B.&1648 ,, &1605 ,, \\
50&H.&3987 ,, &3759 ,, &63&G.&2113 ,, &2045 ,, \\
51&S.&735 ,, &725 ,, &64&M.&4170 ,, &3932 ,, \\
52&P.&2352 ,, &2270 ,, &65&H.&3026 ,, &2894 ,, \\
53&M.&2789 ,, &2676 ,, &66&H.&1916 ,, &1859 ,, \\
55&H.&2813 ,, &2698 ,, &68&M.&4175 ,, &3936 ,, \\\cline{5-8}
57&S.&1798 ,, &1748 ,, &\multicolumn{2}{c}{Summe}&46529 Thl.&44582 Thl.
\end{tabular}
\end{center}
